\chapter{Calculators}
\label{ch:calculators}

\newthought{Use \yamlkey{Calculators}} when your workflow has to drive an external program---a
compiled \textsc{hep} package, a shell wrapper, or any file-based toolchain. \Jarvis\ treats each
program as a black box: it writes the input, runs the program, and reads the output. This chapter
shows the fixed shape every calculator follows.

\section{The black-box idea}
\label{sec:calc-blackbox}

You do not need to understand a tool's internals---only three things:

\begin{enumerate}
  \item what input format it expects (\textsc{json}, \textsc{slha}, plain text);
  \item how to run it (a command line);
  \item what output it produces and which values to read.
\end{enumerate}

\noindent That is exactly what a calculator module declares. Any \textsc{hep} tool---spectrum
generator, dark-matter code, collider recaster---fits the same pattern.

\section{The five phases}
\label{sec:calc-phases}

Every calculator module is described by the same five phases, in this order:

\begin{description}[style=nextline, leftmargin=1.2em]
  \item[\yamlkey{installation}] prepare the program once per instance (copy, unpack, compile).
  \item[\yamlkey{initialization}] reset the working directory before each sample (copy a fresh
        template, delete stale outputs).
  \item[\yamlkey{execution.input}] write this sample's input file(s) (Chapter~\ref{ch:io-types}).
  \item[\yamlkey{execution.commands}] run the program.
  \item[\yamlkey{execution.output}] read observables back out (Chapter~\ref{ch:io-types}).
\end{description}

\section{The module keys}
\label{sec:calc-keys}

Calculator modules live under \yamlkey{Calculators.Modules}; a top-level
\yamlkey{make\_paraller} sets the worker count (Table~\ref{tab:calc-keys}).

\begin{table}[ht]
  \footnotesize
  \begin{tabular}{@{}lp{0.62\linewidth}@{}}
    \toprule
    \textbf{Key} & \textbf{Meaning} \\
    \midrule
    \yamlkey{name}            & module name, used in the workflow graph and logs \\
    \yamlkey{required\_modules} & modules that must finish before this one runs \\
    \yamlkey{clone\_shadow}   & give each sample its own working copy (recommended \code{true}) \\
    \yamlkey{path}            & the program's runtime directory \\
    \yamlkey{source}          & where the program's source/payload lives \\
    \yamlkey{installation}    & commands to build/prepare the program (once) \\
    \yamlkey{initialization}  & commands to reset state before each sample \\
    \yamlkey{execution}       & \yamlkey{path}, \yamlkey{commands}, \yamlkey{input}, \yamlkey{output} \\
    \yamlkey{timeout}         & optional total seconds for one \yamlkey{execution} (see below) \\
    \bottomrule
  \end{tabular}
  \caption{Calculator module keys.}
  \label{tab:calc-keys}
\end{table}

\begin{jwarn}
The worker-count key is spelled \yamlkey{make\_paraller} --- a deliberate legacy spelling. Keep
it exactly; \yamlkey{make\_parallel} will not be recognised. You can also reuse it inside build
commands, e.g. \code{make -j\$\{Calculators:make\_paraller\}}.
\end{jwarn}

\section{Per-sample isolation: \texttt{clone\_shadow} and \token{@PackID}}
\label{sec:calc-shadow}

Most external tools write into their own directory, so running several samples in the same
directory would let them clobber each other. With \yamlkey{clone\_shadow: true}, each sample runs
in its own \emph{shadow} copy, and the \token{@PackID} token in a path expands to that copy's
unique id:

\begin{yamlwin}
path: "&J/calculators/runtime/program/MyCode/@PackID"
\end{yamlwin}

\noindent Use \yamlkey{clone\_shadow: true} unless you have a strong reason not to.

\section{A worked example}
\label{sec:calc-example}

Here is a complete module that drives a spectrum generator: it copies and builds the program,
prepares a template input, writes two parameters into it, runs the tool, and reads two masses
back out.

\begin{yamlwin}[Calculators block]
Calculators:
  make_paraller: 4
  path: "&J/calculators/runtime/program"
  Modules:
    - name: Spectrum
      required_modules: [Parameters]
      clone_shadow: true
      path: "&J/calculators/runtime/program/Spectrum/@PackID"
      source: "&J/deps/program/spectrum"
      installation:
        - "cp ${source}/spectrum.tar.gz ${path}"
        - "cd ${path}"
        - "tar -zxvf spectrum.tar.gz"
        - "make -j${Calculators:make_paraller}"
      initialization:
        - "cp ${source}/template.in ${path}/run.in"
        - "rm -f ${path}/run.out"
      timeout: 120
      execution:
        path: "&J/calculators/runtime/program/Spectrum/@PackID"
        commands:
          - "./run run.in"
        input:
          - name: run_input
            path: "&J/calculators/runtime/program/Spectrum/@PackID/run.in"
            type: "SLHA"
            save: true
            actions:
              - type: "Replace"
                variables:
                  - {name: "M1", placeholder: ">>>M1<<<"}
                  - {name: "M2", placeholder: ">>>M2<<<"}
        output:
          - name: run_output
            path: "&J/calculators/runtime/program/Spectrum/@PackID/run.out"
            type: "xSLHA"
            save: true
            variables:
              - {name: mh1, block: MASS, entry: 25}
              - {name: mN1, block: MASS, entry: 1000022}
\end{yamlwin}

\noindent The output names (\yamlkey{mh1}, \yamlkey{mN1}) are now available to the likelihood and
to any downstream module---this is where an external package becomes useful to the rest of the
scan.

\begin{jtip}
Onboard a new package with a \emph{fixed point} first: set \yamlkey{Sampling.Method} to
\sampler{CSV} (Chapter~\ref{ch:basic-samplers}) and run one or two known points. This validates
install, input, run, and output separately---before you add sampler complexity. Once it works,
drop the same \yamlkey{Calculators} block into your real scan.
\end{jtip}

\section{The full lifecycle}
\label{sec:calc-lifecycle}

For each sample, \Jarvis\ runs the fixed sequence:

\begin{enumerate}
  \item create the sample and allocate a calculator instance directory;
  \item \yamlkey{installation} (once per instance): copy, unpack, build;
  \item \yamlkey{initialization}: reset per-sample input/output state;
  \item \yamlkey{execution.input}: write this sample's input file(s);
  \item \yamlkey{execution.commands}: run the program;
  \item \yamlkey{execution.output}: read the result into observables;
  \item the \yamlkey{LogLikelihood} scores the point, and \Jarvis\ writes one database row.
\end{enumerate}

\section{The execution timeout}
\label{sec:calc-timeout}

\yamlkey{timeout} is an optional safety cut: a total number of seconds for one sample's
\yamlkey{execution} section (input writes, commands, and output reads). It starts after
\yamlkey{initialization}. If omitted, there is no calculator-level limit. Use it to stop a sample
that hangs from stalling the whole scan.

\begin{jnote}
A separate, lower-level per-command limit lives under \yamlkey{Runtime.Subprocess}
(Chapter~\ref{ch:runtime-block}). When both are set, the effective command limit is bounded by
the remaining \yamlkey{timeout} budget.
\end{jnote}

\section{Summary}
\label{sec:calc-summary}

A calculator wraps a program as five phases---install, initialize, write input, run, read
output---with per-sample isolation via \yamlkey{clone\_shadow} and \token{@PackID}. The next
chapter, Chapter~\ref{ch:io-types}, details the input and output file formats those phases use.
